\subsection{Authority Claims}

From Corollary~\ref{cor:complete-discovery}, runtime federation in discovery-based
querying is impossible until every source is discovered. With source counts in the
thousands~\todo{cite here}, this is infeasible.

However, while classical query semantics require all data to be considered in query execution,
this is not always an intuitive requirement. 
Consider a federation of product catalogues, in which a manufacturer publishes the
specifications of the products it makes while resellers and comparison services
republish copies of them. A query over a product's specification is answered from the
manufacturer alone. However, if any copies of, different, or wrong specifications exist 
in the federation, they are considered equally `true'.


We will work towards a query execution model that resolves this observation while
also allowing the delegation of queries \emph{during} query execution.
First we introduce the concept of a \emph{composite source},
an aggregating view over sources in $\Du$ which contains no data of its own.

\begin{definition}[Composite source]
\label{def:composite}
  Fix a setting \\ $W = \langle \Du, \dataf, \Dd_0, \disc \rangle$. A source $c \in \Du$ is
  \emph{composite} when it holds no data of its own and is instead given by a set
  $\Dd_c \subseteq \Du$ of sources it aggregates. It contains all of their data,
  \[
    \dataf(c) \;=\; \bigcup_{d \in \Dd_c} \dataf(d),
  \]
  and a request to $c$ is evaluated over all of it. Sources that are not composite are
  \emph{base} sources. We require the aggregation relation to be well-founded 
  such that no source aggregates itself either directly or through a chain. 
  Thus, when unfolding $\dataf(c)$ it is expressed
  in base sources.
\end{definition}

In the case $\Dd_c = \{d\}$ with $d$ a base source, $\dataf(c) = \dataf(d)$. Then,
$c$ is equivalent to $d$.

\todo{Add this to after claim is defined} 
and interchangeable when $c$ claims (Def.~\ref{def:claim}) nothing. 

% as far as its data goes, $c$ \emph{is} $d$. A degenerate composite that claims nothing
% (Definition~\ref{def:claim}) is interchangeable with $d$ under every definition below, so
% in proofs we may treat a base source as a degenerate composite that claims nothing, and a
% single document that does claim as a degenerate composite that does.


Next we introduce \emph{claims}, which declare a subset of data the composite source
aggregates over that matches a certain criterion. 

\begin{definition}[Claim]
\label{def:claim}
A composite source $c$ may \emph{declare} a claim. A claim is a \emph{triple-wise}
criterion $\sigma_c : \Tt \to \{0, 1\}$, an indicator function mapping triples to 
claim membership.

% i.e.\ $\sigma_c(G) = \{\, t \in G \mid \sigma_c(\{t\}) = \{t\} \,\}$, and $c$
% \emph{claims} $\tau$ when $\sigma_c(\{\tau\}) = \{\tau\}$: the triples $\sigma_c$ selects
% are to be taken from $c$ and from no other source. Write $\mathcal{C} \subseteq \Du$ for
% the composite sources that declare.
\end{definition}


A claim is published with $c$ and thus observed once $c$ is retrieved. 
The core idea on which our query execution model builds is the \emph{principle of
exclusive claims}: queries are answered as if every claim were exclusive, that is, as if
no source other than $c$ held a triple claimed by $c$.
Like the closed-world assumption, this does not state how the federation environment
is; it fixes how we answer queries over it. The data need not satisfy it, and on the Web
it often does not. Definition~\ref{def:claim-semantics} makes the principle precise by
fixing the answer on any data, and the engine enforces that answer, which we discuss in
Section~\todo{?}

% Where the data violate it, we answer over the largest part of
% the data that satisfies it:
% \[
%   G^{\alpha}_A \;=\; \bigcup_{d \in A}
%   \bigl\{\, t \in \dataf(d) \;\bigm|\; d \text{ claims } t
%   \ \text{ or no } c \in \Dd \text{ claims } t \,\bigr\}.
% \]

% \begin{proposition}
% Under Assumption~\ref{as:disjoint}, $G^{\alpha}_{\Dd}$ is obtained by removing from each
% source exactly the triples Assumption~\ref{as:exclusive} forbids it to hold, and no
% others.
% \end{proposition}

Claims indicate that a composite source is the only source for answers for triples that matched
the claim indicator function. Engines can then reason about this claim to determine
for what data delegating BGPs to composite sources declaring a claim is correct.
As an example, in the product catalogue federation,
the manufacturer declares that statements about its products are to be taken from it
alone, so that copies held by resellers are disregarded.

Nothing so far prevents a source from claiming data it should not be allowed to. 
On the Web, any source could claim authority over, for example, DBpedia without being
affiliated with it, and thereby drop every DBpedia result held elsewhere. We
therefore bound each claim by what its source is entitled to.

\begin{definition}[Admissibility, effective claim]
\label{def:admissible}
An \emph{admissibility function} assigns every composite $c$ a criterion
$\rho_c : \Tt \to \{0,1\}$, the triples $c$ is entitled to claim. It is part of the
setting, not of any declaration, and depends only on the triple and the identity of $c$.
The \emph{effective claim} of $c$ is
\[
  \alpha_c(t) \;=\; \sigma_c(t) \cdot \rho_c(t),
\]
with $\alpha_d \equiv 0$ for base sources and for composites that declare nothing. A
source $d$ \emph{claims} $t$ when $\alpha_d(t) = 1$, and the principle of exclusive claims
is read with effective claims.
\end{definition}

An illegitimate claim thus shrinks to its admissible part rather than being rejected
whole: a composite claiming both DBpedia and its own namespace keeps the latter.

\begin{definition}[Claim semantics]
\label{def:claim-semantics}
A triple $t \in \dataf(d)$ is \emph{admitted from} $d$ if
\begin{equation}
  \label{eq:admitted}
  \alpha_d(t) = 1 \quad\text{or}\quad \alpha_c(t) = 0 \text{ for all } c \in \Dd \setminus \{d\},
\end{equation}
and for $A \subseteq \Dd$ the \emph{claim-respecting graph} is
\[
  G^{\alpha}_A \;=\; \bigcup_{d \in A}
  \{\, t \in \dataf(d) \mid t \text{ is admitted from } d \,\}.
\]
The
\emph{claim semantics} of $B$ is $\ev{B}{G^{\alpha}_{\Dd}}$. It replaces the reference
answer of Definition~\ref{def:reach}, and correctness of a delegation
(Definition~\ref{def:delegation}) is measured with $G^{\alpha}_A$ in place of $G_A$.
\end{definition}

Thus, any claimed triple only occurs within the source that claims it, while unclaimed
triples are left in the data as-is. If no source claims
anything, $G^{\alpha}_A = G_A$ and the claim semantics is the reference answer of
Definition~\ref{def:reach}.

\todo{Possibly tie to literature here about subwebs and relational?}

We keep admissibility abstract until Section~\ref{sec:namespace-claims}, where we
instantiate it by namespaces; the results in between hold for any admissibility function.

\begin{assumption}[Disjoint claims]
\label{as:disjoint}
No triple is claimed by two sources: for $c \neq c'$ and every $t \in \Tt$,
$\alpha_c(t) \cdot \alpha_{c'}(t) = 0$.
\end{assumption}

Section~\ref{sec:namespace-claims} gives a sufficient condition for namespace claims. We
discuss relaxing the assumption in Section~\ref{sec:overlap}.~\todo{What happens if we
relax this?}



% \emph{Authority} is the question of \emph{who may speak} about a term. Following the
% Linked Data principles, the document obtained by dereferencing an IRI is the
% distinguished source of statements about that IRI. \citet{hogan2009authoritative}
% made this operational to prevent \emph{ontology hijacking}, and
% \citet{harth2012completeness} generalised it into a lattice of \emph{authority types}
% that record whether a triple was retrieved from the document of its subject, its
% predicate, or its object. Authority is a \emph{filter}: it discards data, and it is
% chosen by the consumer.

% \emph{Completeness} is the question of \emph{what is guaranteed to be present}. In the
% relational setting this is Motro's split of a source into an available and an ideal
% instance~\cite{motro1989integrity}, and Levy's local completeness
% statements~\cite{levy1996complete,razniewski2011completeness};
% \citet{darari2013completeness,darari2018management} transported it to RDF as
% statements $\mathrm{Compl}(P_1 \mid P_2)$ asserting that a source holds every ideal
% answer to $P_1$ under condition $P_2$. Completeness is an \emph{obligation}: it adds
% nothing to the data, and it is undertaken by the publisher.

% Neither ingredient alone licenses what we need. Authority alone does not let us skip a
% source; completeness alone does not let us disregard the statements other sources make
% about the same subjects. Their conjunction does, and we make that conjunction the one
% assumption on which the rest of the paper rests.

% \subsubsection{Data model}

% We use the standard model of a Web of Linked Data~\cite{hartig2012foundations},
% refined with partial documents as in~\cite{bogaerts2021subweb}. Let $\mathcal{U}$,
% $\mathcal{B}$ and $\mathcal{L}$ be the pairwise disjoint sets of IRIs, blank nodes and
% literals, and $\Tt = (\mathcal{U} \cup \mathcal{B}) \times \mathcal{U} \times
% (\mathcal{U} \cup \mathcal{B} \cup \mathcal{L})$ the set of triples. A \emph{Web of
% Linked Data} (WoLD) is a tuple $W = \langle D, \mathsf{data}, \mathsf{adoc} \rangle$
% with $D$ a countable set of documents, $\mathsf{data} : D \to 2^{\Tt}$ mapping
% each document to a finite graph, and $\mathsf{adoc} : \mathcal{U} \rightharpoonup D$
% the partial dereferencing function. We write $\mathsf{AllData}(W) = \bigcup_{d \in D}
% \mathsf{data}(d)$, and $G|_S = \{ (s,p,o) \in G \mid s \in S \}$ for the restriction
% of a graph to subjects drawn from $S \subseteq \mathcal{U}$.

